% ===================== فصل ۱۷ =====================
\chapter{هم‌روندی}
\label{ch:concurrency}

تا اینجا برنامه‌های ما یک کار را در هر لحظه انجام می‌دادند. اما رایانه‌های امروز
چند هسته دارند و می‌توانند چند کار را \emph{هم‌زمان} پیش ببرند. \textbf{هم‌روندی}
(concurrency) هنرِ بهره‌گیری از این توانایی است. در این فصل با نخ‌ها و
هماهنگیِ میانشان آشنا می‌شویم.

\begin{warn}
این فصل از پیشرفته‌ترین مباحثِ کتاب است. اگر تازه‌کارید، می‌توانید فعلاً از آن
بگذرید و پس از تسلط بر فصل‌های پیشین بازگردید. هم‌روندی قدرتمند است، اما نیازمندِ
دقت.
\end{warn}

\section{نخ چیست؟}

یک \textbf{نخ} (thread) مانندِ یک کارگرِ مستقل است که می‌تواند هم‌زمان با
کارگرانِ دیگر کار کند. اگر چهار نخ بسازید و رایانهٔ شما چهار هسته داشته باشد،
هر چهار کار می‌توانند واقعاً هم‌زمان پیش بروند و برنامه تا چهار برابر
سریع‌تر شود.

ماژولِ \code{sync} ابزارهای کار با نخ را فراهم می‌کند. دو تابعِ کلیدی:
\begin{itemize}
  \item \code{spawn(تابع)}: یک نخِ تازه می‌سازد که آن تابع را اجرا می‌کند، و
        یک شناسه برمی‌گرداند.
  \item \code{join(شناسه)}: منتظر می‌ماند تا آن نخ کارش را تمام کند.
\end{itemize}

\section{مسئلهٔ مشترک: داده‌ای که چند نخ به آن دست می‌زنند}

وقتی چند نخ به \emph{یک} داده دست می‌زنند، دردسر آغاز می‌شود. تصور کنید دو نخ
هم‌زمان می‌خواهند به یک شمارنده یکی اضافه کنند؛ ممکن است هر دو مقدارِ کهنه را
بخوانند و یکی از افزایش‌ها گم شود. به این مشکل \textbf{رقابت} (race condition)
می‌گویند.

راهِ حل، \textbf{قفل} (Mutex) است: ابزاری که تضمین می‌کند در هر لحظه فقط
\emph{یک} نخ به داده دست می‌زند. این مثالِ واقعی که جمعِ اعداد را میانِ
چهار نخ تقسیم می‌کند همهٔ این ابزارها را به کار می‌گیرد:

\begin{salamen}
package main
import ("sync" "str" "io")
extern "C":
    func salam_alloc(size u64) void*
    func salam_free(ptr void*)
end
const N = 100000
mut sum i64* = null
mut lock sync.Mutex = sync.NewMutex()
mut wg sync.WaitGroup = sync.NewWaitGroup()
func add_range(lo i32, hi i32):
    mut local i64 = 0
    mut i = lo
    while i < hi:
        local = local + i as i64
        i = i + 1
    end
    mut acc i64* = sum
    sync.Lock(lock)
    acc[0] = acc[0] + local
    sync.Unlock(lock)
    sync.Done(wg)
end
func q1(): add_range(0, N / 4) end
func q2(): add_range(N / 4, N / 2) end
func q3(): add_range(N / 2, N * 3 / 4) end
func q4(): add_range(N * 3 / 4, N) end
func main:
    sum = salam_alloc(8 as u64) as i64*
    sum[0] = 0
    sync.Add(wg, 4 as i64)
    t1 i64 = spawn(q1)
    t2 i64 = spawn(q2)
    t3 i64 = spawn(q3)
    t4 i64 = spawn(q4)
    sync.Wait(wg)
    join(t1)
    join(t2)
    join(t3)
    join(t4)
    io.Print("sum of 0..")
    io.Print(str.FromInt(N as i64))
    io.Print(" = ")
    io.Println(str.FromInt(sum[0]))
    salam_free(sum as void*)
end
\end{salamen}

این برنامه پیچیده به‌نظر می‌رسد، اما ایده‌اش روشن است. بیایید آرام بشکافیمش.

\subsection{تقسیمِ کار}

به‌جای آنکه یک نخ همهٔ اعدادِ ۰ تا ۱۰۰٬۰۰۰ را جمع بزند، کار را به چهار بخش
تقسیم کرده‌ایم. هر یک از توابعِ \code{q1} تا \code{q4} یک‌چهارمِ بازه را جمع
می‌زند. این «تقسیم و غلبه» جانِ کلامِ هم‌روندی است.

\subsection{محافظت با قفل}

هر نخ نخست جمعِ \emph{محلیِ} خود را در \code{local} حساب می‌کند (که هیچ نخِ
دیگری به آن دست نمی‌زند، پس امن است). تنها در پایان، که می‌خواهد نتیجه را به
مجموعِ مشترک بیفزاید، قفل می‌گیرد:

\begin{salamen}
sync.Lock(lock)
acc[0] = acc[0] + local
sync.Unlock(lock)
\end{salamen}

میانِ \code{Lock} و \code{Unlock} فقط یک نخ می‌تواند باشد. پس افزودنِ نتیجه
هرگز دچارِ رقابت نمی‌شود.

\subsection{انتظار با گروهِ انتظار}

\code{WaitGroup} ابزاری است که می‌گوید «منتظرِ پایانِ همهٔ نخ‌ها بمان». با
\code{sync.Add(wg, 4)} گفتیم چهار کار در راه است؛ هر نخ در پایان
\code{sync.Done(wg)} را صدا می‌زند؛ و \code{sync.Wait(wg)} در تابعِ اصلی صبر
می‌کند تا هر چهار «انجام‌شده» اعلام شوند.

\begin{tip}
الگوی استانداردِ هم‌روندیِ ایمن: (۱) کار را به بخش‌های مستقل تقسیم کن، (۲)
بگذار هر نخ روی دادهٔ \emph{محلیِ} خودش کار کند، (۳) فقط در پایان و با
\emph{قفل}، نتیجه‌ها را با هم بیامیز، (۴) با \code{Wait} و \code{join} منتظرِ
پایانِ همه بمان. اگر این الگو را رعایت کنید، از بیشترِ دردسرهای هم‌روندی در امان
می‌مانید.
\end{tip}

\begin{deep}[نگاهی به درون: چرا جمعِ محلی؟]
اگر هر نخ \emph{در هر گام} قفل می‌گرفت تا به مجموعِ مشترک اضافه کند، نخ‌ها
بیشترِ وقتشان را پشتِ قفل منتظر می‌ماندند و هم‌روندی هیچ سودی نداشت حتی شاید
کندتر می‌شد! ترفندِ «جمعِ محلی، سپس یک‌بار آمیختن» تعدادِ دفعاتِ قفل‌گیری را از
صدها‌هزار به فقط چهار کاهش می‌دهد. این یک اصلِ مهم است: \emph{قفل را تا جای
ممکن کوتاه و کم نگه دارید.}
\end{deep}

\section{خطرهای هم‌روندی}

هم‌روندی قدرتمند است اما تله‌هایی دارد که باید بشناسید:
\begin{itemize}
  \item \textbf{رقابت:} دو نخ بی‌محافظت به یک داده دست می‌زنند. \emph{درمان:}
        قفل.
  \item \textbf{بن‌بست (deadlock):} دو نخ هر کدام منتظرِ قفلی هستند که دیگری
        گرفته، و هیچ‌کدام پیش نمی‌رود. \emph{درمان:} همیشه قفل‌ها را به ترتیبِ
        یکسان بگیرید.
  \item \textbf{فراموشیِ آزادسازی:} اگر \code{Unlock} را فراموش کنید، نخ‌های
        دیگر تا ابد منتظر می‌مانند. \emph{درمان:} \code{Lock} و \code{Unlock}
        را همیشه جفت بنویسید.
\end{itemize}

\section{خلاصهٔ فصل}

\begin{itemize}
  \item نخ کارگری مستقل است که هم‌زمان با دیگران کار می‌کند.
  \item با \code{spawn} نخ می‌سازیم و با \code{join} منتظرش می‌مانیم.
  \item قفل (\code{Mutex}) تضمین می‌کند فقط یک نخ به داده دست می‌زند.
  \item \code{WaitGroup} برای انتظارِ پایانِ همهٔ نخ‌هاست.
  \item بهترین الگو: کارِ مستقل، دادهٔ محلی، آمیختنِ کوتاه با قفل.
  \item مراقبِ رقابت و بن‌بست باشید.
\end{itemize}

\section{تمرین‌ها}

\exercise برنامه‌ای بنویسید که دو نخ بسازد و هر کدام پیامِ متفاوتی چاپ کند، و در
تابعِ اصلی منتظرِ هر دو بماند.

\exercise جمعِ اعدادِ ۱ تا ۱۰۰۰۰ را یک بار با یک نخ و یک بار با چهار نخ حساب
کنید و نتیجه‌ها را مقایسه کنید.

\exercise توضیح دهید چرا در مثالِ این فصل، استفاده از \code{local} به‌جای
افزودنِ مستقیم به \code{sum} مهم است.

\exercise یک سناریوی بن‌بست را روی کاغذ شرح دهید و بگویید چگونه می‌توان از آن
پرهیز کرد.

\exercise بررسی کنید اگر در مثالِ این فصل \code{sync.Lock} و \code{sync.Unlock}
را حذف کنیم، چه مشکلی ممکن است پیش بیاید.
